\section{Limitations}
\label{sec:limitations}

\paragraph{Common-law doctrines are not modelled.} Yuho captures the
statute as drafted; it does not model judicial gloss. Doctrines like the
\emph{thin skull} rule or \emph{novus actus interveniens} live in case
law and shape how elements (especially causation) apply. Our \texttt{caselaw}
block records the existence of judicial interpretation but does not lift
the doctrine into the type system.

\paragraph{Single jurisdiction.} The encoding is Singapore-only. The
Anglo-Indian heritage of the Singapore Penal Code makes the grammar
plausibly portable to India, Pakistan, Brunei, and (with adjustments)
Malaysia, but we have not verified this empirically.

\paragraph{Shallow precedent integration.} Case law is referenced by name
+ citation + holding string; the holding is not parsed into structured
elements. A holding like ``mens rea may be inferred from the actus reus
when the act is inherently dangerous'' would, in a richer system, modify
the burden qualifier of a mens-rea element. Yuho does not yet do this.

\paragraph{No live SSO diffing.} The canonical scrape lives in
\texttt{library/penal\_code/\_raw/act.json} as a snapshot. Amendments
published by AGC are not picked up automatically. A scheduled re-scrape
+ diff pass is on the roadmap.

\paragraph{L3 reviewer is the encoder.} The 11-point checklist is
mechanical, but the same author wrote the encodings and stamped them.
External review by a Singapore-qualified lawyer is the natural next step
for any encoding intended for use beyond proof of concept.

\paragraph{Verification scope.} Z3 and Alloy hookups currently target
exception-priority well-formedness and bounded element-combination
enumeration. Properties that require quantification over an unbounded
number of historical events (e.g.\ ``no defendant has ever been convicted
under both s299 and s300 for the same act'') are out of scope.
